\section{\texorpdfstring{\texttt{def}, \texttt{let}, implicit arguments}{def, let, implicit arguments}}\label{sec:01-basics:02-def-let-implicit:1}

Section 1 established that every term has a type, and that Lean checks this before running anything. That fact is useless without a way to \emph{name} terms so they can be reused, and a way to write one definition that works for many types at once rather than one copy per type. Those are exactly the two gaps this section closes. \texttt{def}/\texttt{let} handle naming (global and local, respectively), and implicit arguments allow writing a single definition, like an identity function, that is correct for every type simultaneously, instead of retyped per type.

Consider three short definitions, each introducing one of these pieces in turn. Every token in them is doing something specific; each is examined in full below.

\begin{lstlisting}
def double (n : Nat) : Nat := n * 2
\end{lstlisting}

\textbf{\texttt{def double (n : Nat) : Nat := n * 2}}

\begin{itemize}
\tightlist
\item
  \texttt{def} is the keyword that introduces a new global definition into the environment, a name, permanently available afterward, standing for a fixed term. This is different from a hypothesis or a local variable. Once elaborated, \texttt{double} is a completely ordinary constant, referable anywhere later in the file (or, once the file is imported, anywhere else). Compare this with \texttt{example} from Chapter 4, which elaborates a term but does \emph{not} bind it to a name. \texttt{def} is for the case where the definition is meant to be reused.
\item
  \texttt{double} is the name being bound. Lean enforces no special naming convention, but lowerCamelCase for \texttt{def}s and UpperCamelCase for types is the near-universal community style, which this book follows.
\item
  \texttt{(n : Nat)} is an explicit argument named \texttt{n}, of type \texttt{Nat}. ``Explicit'' means a caller must supply it positionally. \texttt{double 5} provides \texttt{5} for \texttt{n} directly. This is what makes \texttt{double} a function rather than a plain value. Everything after \texttt{def double} up to the first bare \texttt{:=} (if there are argument binders) is the function parameter list, and the actual type of \texttt{double} ends up being \texttt{Nat → Nat}, exactly as if it had been written \texttt{def double : Nat → Nat := fun n => n * 2} instead. The two forms elaborate to the same term. The \texttt{(n : Nat)} binder form is simply the standard, more readable surface syntax for ``a function with named parameters.''
\item
  \texttt{: Nat} (the second one, right before \texttt{:=}) is the declared return type. This is not optional filler. Lean uses it to \emph{check} the body against a known expected type while elaborating, rather than only inferring a type afterward. If the type of the body did not match, an error would occur at this point, not somewhere downstream.
\item
  \texttt{:=} separates the declaration signature (name, arguments, return type) from the definition (the actual term). Read it as ``is defined to be.''
\item
  \texttt{n * 2} is the body. At this point in elaboration, \texttt{n} is in scope with type \texttt{Nat} (introduced by the \texttt{(n : Nat)} binder above), so \texttt{n * 2} is \texttt{Nat} multiplication applied to \texttt{n} and the numeral \texttt{2}. That numeral itself elaborates to a \texttt{Nat}, because that is what the type of the second argument of \texttt{Nat.mul} forces it to be. Numeral elaboration is guided by the expected type, another case of Lean checking against context instead of guessing.
\end{itemize}

A word of warning about the \texttt{→} that appeared two bullets up (\texttt{Nat → Nat}, the type of \texttt{double}). This same arrow reappears twice more later in this book, meaning something that looks different each time. Once \texttt{Prop} exists (\hyperref[sec:04-propositions-and-proofs:04-implication:1]{Chapter 4}), \texttt{P → Q} for propositions \texttt{P}, \texttt{Q} reads as logical implication, not ``takes a \texttt{P}-value, returns a \texttt{Q}-value.'' And later, \hyperref[sec:02-terminology-and-coc:02-pi-sigma-and-coc:1]{Chapter 2, Section 2} names the fully general pattern \texttt{→} is secretly always an instance of, the Π-type, where the \emph{codomain} is allowed to depend on which argument was given, \texttt{Nat → Nat} being only the special case where it happens not to. Nothing about \texttt{double} above needs any of this yet. It is flagged only so the same symbol showing up wearing two more meanings later does not read as three unrelated coincidences.

\begin{mathreading}

\texttt{def double (n : Nat) : Nat := n * 2} is nothing more than the ordinary mathematical definition

\end{mathreading}

\[
\mathrm{double} : \mathbb{N} \to \mathbb{N}, \qquad \mathrm{double}(n) = 2n,
\]

with the signature \(\mathrm{double} : \mathbb{N} \to \mathbb{N}\) split across \texttt{def double (n : Nat) : Nat} (domain and codomain, spelled out argument-by-argument rather than as a single arrow), and the equation \(\mathrm{double}(n) = 2n\) becomes \texttt{:= n * 2}. There is no real difference between writing the domain as one arrow \texttt{Nat → Nat} or as a named binder \texttt{(n : Nat) : Nat}, exactly as \(f : A \to B,\ f(a) = \ldots\) and \(f = (a \mapsto \ldots) : A \to B\) describe the same function.

\begin{lstlisting}
def average (a b : Nat) : Nat :=
  let sum := a + b
  sum / 2
\end{lstlisting}

\textbf{\texttt{def average (a b : Nat) : Nat := let sum := a + b; sum /\allowbreak{} 2}}

\begin{itemize}
\tightlist
\item
  \texttt{(a b : Nat)} is two explicit arguments, both of type \texttt{Nat}, written with a single shared type annotation. This is pure surface-syntax sugar for \texttt{(a : Nat) (b : Nat)}. Lean expands it identically either way. This shorthand is standard style whenever several consecutive parameters share a type, and costs nothing.
\item
  \texttt{let sum := a + b} introduces a \emph{local} definition, visible only in the rest of this particular body, as opposed to the \emph{global} kind introduced by \texttt{def}. \texttt{let} does not need (though it can take) a type annotation, since the type of \texttt{a + b} is already fully determined by the types of \texttt{a} and \texttt{b}, so Lean infers it. Operationally, \texttt{let sum := a + b; sum /\allowbreak{} 2} means exactly the same thing as substituting \texttt{a + b} for every occurrence of \texttt{sum} in \texttt{sum /\allowbreak{} 2}. A \texttt{let} is definitionally transparent, so \texttt{average a b} and \texttt{(a + b) /\allowbreak{} 2} elaborate to the same normal form. The reason to write it as a \texttt{let} anyway, rather than inlining \texttt{(a + b) /\allowbreak{} 2} directly, is purely for the human reader. Naming an intermediate quantity documents what it means, and in a longer proof or definition, it prevents repeating a nontrivial subexpression (and therefore repeating a mistake in it) in several places.
\item
  The line break between \texttt{let sum := a + b} and \texttt{sum /\allowbreak{} 2} is whitespace, not two separate statements needing a semicolon. The parser used by Lean uses indentation-sensitive layout for \texttt{let}-chains, the same way it does for \texttt{by}-blocks in tactic mode. \texttt{sum /\allowbreak{} 2} is the body of the \texttt{let}. The whole two-line construct \texttt{let sum := a + b; sum /\allowbreak{} 2} is itself one term, which is then what the \texttt{:=} of \texttt{average} binds to.
\item
  \texttt{sum /\allowbreak{} 2} is \texttt{Nat} division, which in Lean is \emph{truncating}. \texttt{average 4 10} computes \texttt{sum = 14}, then \texttt{14 /\allowbreak{} 2 = 7} exactly. But \texttt{average 1 2} would compute \texttt{sum = 3}, then \texttt{3 /\allowbreak{} 2 = 1} (rounded down, since \texttt{Nat} has no fractions). This should be noted before relying on this \texttt{average} for anything where the rounding matters.
\end{itemize}

\begin{mathreading}

The \texttt{let} is exactly a ``let \(s := a + b\) in \(\ldots\)'' clause of the kind used constantly in written proofs to name an intermediate quantity.

\end{mathreading}

\[
\mathrm{average}(a,b) = \big\lfloor \tfrac{s}{2} \big\rfloor
\text{ where } s := a + b,
\]

One caveat is that this \texttt{average} computes \(\lfloor s/2 \rfloor\) (floor division) rather than the true rational average \(s/2\), since division on \texttt{Nat} is truncating. This gap is easy to overlook in ordinary mathematical prose, but Lean forces it to be confronted explicitly. There is no coercion to \(\mathbb{Q}\) happening for free.

\begin{lstlisting}
def identity {(*$\alpha$*) : Type} (x : (*$\alpha$*)) : (*$\alpha$*) := x
\end{lstlisting}

\textbf{\texttt{def identity \{α : Type\} (x : α) : α := x}}

\begin{itemize}
\tightlist
\item
  \texttt{\{α : Type\}} is an \textbf{implicit} argument, marked by curly braces instead of parentheses. The name \texttt{α} (conventionally a Greek letter for a type variable; again pure convention, \texttt{T} or \texttt{A} would work just as well) has type \texttt{Type}, meaning this argument is itself a \emph{type}, not a value of some fixed type. \texttt{identity} is thus \textbf{polymorphic}. It works uniformly for every choice of \texttt{α}.
\item
  The crucial difference from \texttt{(n : Nat)} above is that \texttt{\{α : Type\}} is not supplied positionally at call sites. Writing \texttt{identity 5} does \emph{not} mean ``pass \texttt{5} as \texttt{α}.'' Lean instead \textbf{elaborates} (infers) \texttt{α} by unification, working backward from the type of the explicit argument actually supplied. In \texttt{identity 5}, Lean sees that \texttt{5 : Nat} is being passed where an \texttt{x : α} is expected, unifies \texttt{α := Nat}, and only then checks the rest. An implicit argument can still be supplied explicitly when overriding inference is necessary, with \texttt{@identity Nat 5}. The \texttt{@} prefix means ``no more auto-inference; every argument, implicit or not, is given by hand.'' This escape hatch matters mainly for debugging elaboration failures, not everyday use.
\item
  The reason to make \texttt{α} implicit rather than explicit here is that the value of \texttt{α} is determined by the type of \texttt{x} at every call site, so requiring the caller to type it out (as in \texttt{identity Nat 5}) would be pure noise. Lean already has enough information without being told. The general rule of thumb, used throughout Mathlib and this book, is to mark an argument implicit exactly when its value is always recoverable from the \emph{other} arguments or from the expected return type, and to keep it explicit when it genuinely varies independently, and a reader benefits from seeing it written at the call site.
\item
  \texttt{: α := x} gives the return type \texttt{α} itself (the same type variable bound above, now in scope for the rest of the signature and body), and the body is simply \texttt{x}, the argument unchanged. This is the identity function at every type simultaneously, one \texttt{def}, rather than one per type, which is exactly what the \texttt{\{α : Type\}} parameter provides.
\end{itemize}

\begin{mathreading}

\texttt{identity \{α : Type\} (x : α) : α := x} is the family \(\{\,\mathrm{id}_A : A \to A\,\}_{A \in \mathbf{Type}}\) indexed over \emph{every} type \(A\) at once. It is precisely the assignment sending each object \(A\) of the category to its identity morphism \(\mathrm{id}_A\), packaged as a single polymorphic definition rather than one definition per \(A\). The implicit argument \texttt{\{α : Type\}} is what makes this a \emph{statement about all \(A\) simultaneously} rather than a single fixed function. In category theory one would never write ``\(\mathrm{id}_{\mathbb{Z}}\), \(\mathrm{id}_{\mathbb{R}}\), \ldots{}'' one at a time either, but would instead say ``for every object \(A\), there is an identity morphism \(\mathrm{id}_A\),'' which is exactly what the universally-quantified, implicitly-inferred \texttt{\{α : Type\}} expresses.

\end{mathreading}

\textbf{Two more binder styles, and one more kind of \texttt{def}, not needed yet.} \texttt{()} and \texttt{\{\}} are not the whole story, only the two pieces needed so far. Briefly, so meeting these elsewhere later is not a surprise:

\begin{itemize}
\tightlist
\item
  A \textbf{third} kind of argument, written with square brackets, \texttt{[x : C]}, is solved neither positionally like \texttt{()} nor by unification like \texttt{\{\}}, but by a search through registered instances of \texttt{C}, called \textbf{typeclass resolution}. This book delays it deliberately, until \hyperref[sec:06-rigor-check:01-structure-vs-class:1]{Chapter 6}, where the \texttt{Group}/\texttt{class}/\texttt{instance} machinery that makes it useful is built up from \texttt{structure} first, so the automation is understood rather than taken on faith.
\item
  A \textbf{fourth}, rarer kind, \textbf{strict implicit}, written \texttt{\{\{x : α\}\}} (Lean also accepts a Unicode spelling of the same double-brace pair, seen occasionally in Mathlib source), behaves like \texttt{\{\}} except that Lean defers solving it until an \emph{explicit} argument after it is actually supplied. It is a Mathlib idiom for keeping partially applied functions well behaved, and does not appear anywhere in the own code of this book. It is named here only so it is recognizable, not mysterious, if encountered while reading Mathlib source directly.
\item
  \texttt{def} is not the only way to introduce a definition either. \texttt{abbrev} and \texttt{opaque} exist alongside it, and differ not in what they let you write, but in how transparent the result is to the own equality checker of Lean, whether \texttt{unfold}ing it is ever needed, or ever even possible. \hyperref[sec:05-tactics:04-more-tactics:1]{Chapter 5}, once \texttt{unfold} itself is on the table, is where this is worth actually seeing rather than taking on faith.
\end{itemize}

None of this needs remembering yet. Only \texttt{()} and \texttt{\{\}} are needed for everything through the end of this chapter.

The implicit \texttt{\{α : Type\}} argument of \texttt{identity} is polymorphic \emph{over which type is plugged in}, but its return type, \texttt{α}, never changes shape once \texttt{α} is fixed. An \texttt{identity} for \texttt{Nat} returns a \texttt{Nat}, full stop, no matter which \texttt{Nat} was given. The next section asks the sharper question this naturally raises. Can the \emph{type itself} depend not just on which type was chosen, but on the specific \emph{value} of an argument, changing from one call to the next even at a single fixed type? That is a strictly new capability, not a variation on implicit arguments, and it is what makes Lean a genuine proof assistant rather than an ordinary typed language.
